\input{../output/tables/q4_numbers.tex}

\subsection{价格结构与因果预测}

附件四价格范围为0.008--1.794元/kWh，全年均值与附件一相同，平均日内轮廓相关系数为1.0；主要变化体现在日级水平和日内偏差（图~\ref{fig:q4-prices}）。逐时价格与净负荷的相关系数为\QFourPriceNetCorr，逐时段回归$R^2$为\QFourPriceRSquared。净负荷预测误差最大的20\%时段，电价为均值的\QFourDeficitPriceRatio 倍，提示补购风险与高价可能叠加。

据此，以最近两个同星期日的价格轮廓与最近56日逐时段净负荷回归等权预测价格，回归输入采用本地净负荷预测，两个提前期共用已完成日期。组合由一月8--31日误差检验确定，二月起冻结。6、12、18时再以当日已观测的平均价格偏差修正剩余预测，形成“日前轮廓＋日内偏差”的更新方式。

\input{../output/tables/q4_price.tex}

二月至十二月，融合预测MAE为\QFourMAECombined 元/kWh，低于固定轮廓的\QFourMAEFixed 元/kWh；三次日内修正也降低了剩余时段MAE。预测是否节费，仍需在相同结算规则下通过运行对照判断。

\begin{figure}[htbp]
 \centering\includegraphics[width=\textwidth]{../output/figures/q4_prices.pdf}
 \caption{实时电价的结构与可预测性。a：附件四实时电价按星期与10分钟时段的二月至十二月均值；b：各日平均电价与平均净负荷，周五、周六由价格聚类识别，虚线为最小二乘拟合；c：二月至十二月各因果日前预测方法的当日平均绝对误差，条形从零起。}
 \label{fig:q4-prices}
\end{figure}

\FloatBarrier
\subsection{两类计划调整机制与结算模型}
\label{sec:q4-price-convention}

4-2沿用本地预测并保持0时计划购电量不变，4-3沿用融合光伏预报、三时点调整及其信息链。规划时以预测价格替换相应代理模型中的$p_t$，实际现金费用则按相应时段实价计算：
\begin{equation}
  C_{4\text{-}2}^{\rm cash}=\sum_t p_t(q_t^0+5e_t),
  \label{eq:problem42}
\end{equation}
\begin{equation}
  C_{4\text{-}3}^{\rm cash}
  =\sum_t p_t(q_t^0+1.5\delta_t^+-0.5\delta_t^-+5e_t).
  \label{eq:problem43}
\end{equation}
实际价格只进入结算和已到达的观测，不提前供规划使用。两族保持设备、执行器和残差构造一致，并分别设置固定轮廓及当日价格预知对照；4-3另设无日内价格修正和仅0时计划。价格预知只替换输入价格，保留代理目标与执行器，因而不构成最优费用界限。

当价格$P$与需求$D$相关时，单时段最优条件为$\mathbb E[P]=5\mathbb E[P\mathbf1_{\{D>q\}}]$，普通分位数不再自动适用。本扩展暂继承0.8分位数与$[0.7,0.9]$调整带以保持机制可比，其结论属于因果点价格预测下的代理调度。

\subsection{运行结果与调整收益}

\textbf{计划固定（4-2）。}现金费用为\QFourTwoTotalCost 万元，其中计划费\QFourTwoPlanCost 万元、紧急费\QFourTwoEmergencyCost 万元；紧急购电\QFourTwoEmergencyEnergy\ kWh，期末电量\QFourTwoTerminal\ kWh，日紧急费用CVaR95为\QFourTwoCvar 元。相对固定电价下的问题二，费用增加\QFourTwoVsQTwoPercent\%，说明在相同平均电价下，价格与购电时序的匹配仍会改变支出。

\textbf{有偿调整（4-3）。}现金费用为\QFourThreeTotalCost 万元，由计划费\QFourThreePlanCost 万元、调整结算\QFourThreeAdjustCost 万元与紧急费\QFourThreeEmergencyCost 万元组成。紧急购电\QFourThreeEmergencyEnergy\ kWh，CVaR95为\QFourThreeCvar 元。相对同族仅0时计划，节省\QFourThreeSavingPlanOnly 元（\QFourThreeSavingPlanOnlyPercent\%），紧急购电由\QFourThreePlanOnlyEmergency 降至\QFourThreeEmergencyEnergy\ kWh。

4-3较4-2节省\QFourThreeVsTwo 元（\QFourThreeVsTwoPercent\%），日均节省95\%区间为$[\QFourCIThreeVsTwoLow,\QFourCIThreeVsTwoHigh]$元。经“融合预报、仅计划”这一中间策略分解，0时融合预报与日内调整分别贡献\QFourValueExternalPlan 元和\QFourValueAdjust 元，均以4-2为分母，对应\QFourValueExternalPlanPercent\%和\QFourValueAdjustPercent\%；日内调整占合计节省的\QFourAdjustShareOfSaving\%，其中包含实际状态反馈的作用。

\input{../output/tables/q4_comparison.tex}

\begin{figure}[htbp]
 \centering\includegraphics[width=\textwidth]{../output/figures/q4_costs.pdf}
 \caption{实时价格下的费用、补购与收益来源。a：两族策略的334天现金费用，括号为相邻变化，橙线为增支、蓝线为节省，横轴不从零起；空心点为价格预知启发式对照，不是最优费用界限。b：对应紧急购电量，从零起。c：4-3相对4-2的节省分解，日内调整包含状态反馈。a的误差线对应相邻差，c的误差线对应合计节省，均为配对日均费用差的七日循环块重采样95\%区间乘以334天。}
 \label{fig:q4-costs}
\end{figure}

\FloatBarrier
\subsection{结算效应、预测收益与风险}

图~\ref{fig:q4-costs}采用逐层替换识别费用来源。保持问题二、问题三的决策不变，仅改用实时价结算，分别增支\QFourTwoSettlementEffect 元和\QFourThreeSettlementEffect 元（\QFourTwoSettlementPercent\%与\QFourThreeSettlementPercent\%）；再让策略响应因果价格预测，分别收回\QFourTwoFixedDiff 元和\QFourThreeFixedDiff 元。4-2的日均预测节费95\%区间为$[\QFourCITwoFixedLow,\QFourCITwoFixedHigh]$元。预测利用了部分价格规律，但尚未抵消结算变化带来的全部增支。

更高的预测精度也不必然减少费用或补购。价格预知对照比两族主策略分别少付\QFourTwoPerfectGap 元和\QFourThreePerfectGap 元，但紧急购电均更高，两族分别有105天和109天的日费用高于因果策略。对4-3，取消日内价格修正仅使费用变化\QFourThreeNoIntradayDiff 元，日均差区间$[\QFourCIThreeNoIntradayLow,\QFourCIThreeNoIntradayHigh]$元包含零；信息链修正的费用差同样不显著，且补购略升。这些结果共同表明，点预测的改善只有在改变关键时段的计划购电量与储能配置时，才会形成实际收益。

指定日期结果见附录表~\ref{tab:q42-purchase}至表~\ref{tab:q43-emergency}。两份工作簿\path{result4-2.xlsx}、\path{result4-3.xlsx}分别沿用问题二、问题三的结果口径。全部\QFourStrategyCount 条全年轨迹、规划价格和交付账单通过独立核验，详见\hyperref[sec:model-validation]{第九章}及附录\ref{app:validation-details}。
